% Options for packages loaded elsewhere
\PassOptionsToPackage{unicode}{hyperref}
\PassOptionsToPackage{hyphens}{url}
\PassOptionsToPackage{dvipsnames,svgnames,x11names}{xcolor}
%
\documentclass[
  letterpaper,
  DIV=11,
  numbers=noendperiod]{scrartcl}

\usepackage{amsmath,amssymb}
\usepackage{iftex}
\ifPDFTeX
  \usepackage[T1]{fontenc}
  \usepackage[utf8]{inputenc}
  \usepackage{textcomp} % provide euro and other symbols
\else % if luatex or xetex
  \usepackage{unicode-math}
  \defaultfontfeatures{Scale=MatchLowercase}
  \defaultfontfeatures[\rmfamily]{Ligatures=TeX,Scale=1}
\fi
\usepackage{lmodern}
\ifPDFTeX\else  
    % xetex/luatex font selection
  \setmainfont[]{Inter}
  \setsansfont[]{Inter}
  \setmathfont[]{Fira Math}
\fi
% Use upquote if available, for straight quotes in verbatim environments
\IfFileExists{upquote.sty}{\usepackage{upquote}}{}
\IfFileExists{microtype.sty}{% use microtype if available
  \usepackage[]{microtype}
  \UseMicrotypeSet[protrusion]{basicmath} % disable protrusion for tt fonts
}{}
\makeatletter
\@ifundefined{KOMAClassName}{% if non-KOMA class
  \IfFileExists{parskip.sty}{%
    \usepackage{parskip}
  }{% else
    \setlength{\parindent}{0pt}
    \setlength{\parskip}{6pt plus 2pt minus 1pt}}
}{% if KOMA class
  \KOMAoptions{parskip=half}}
\makeatother
\usepackage{xcolor}
\usepackage{soul}
\setlength{\emergencystretch}{3em} % prevent overfull lines
\setcounter{secnumdepth}{5}
% Make \paragraph and \subparagraph free-standing
\ifx\paragraph\undefined\else
  \let\oldparagraph\paragraph
  \renewcommand{\paragraph}[1]{\oldparagraph{#1}\mbox{}}
\fi
\ifx\subparagraph\undefined\else
  \let\oldsubparagraph\subparagraph
  \renewcommand{\subparagraph}[1]{\oldsubparagraph{#1}\mbox{}}
\fi


\providecommand{\tightlist}{%
  \setlength{\itemsep}{0pt}\setlength{\parskip}{0pt}}\usepackage{longtable,booktabs,array}
\usepackage{calc} % for calculating minipage widths
% Correct order of tables after \paragraph or \subparagraph
\usepackage{etoolbox}
\makeatletter
\patchcmd\longtable{\par}{\if@noskipsec\mbox{}\fi\par}{}{}
\makeatother
% Allow footnotes in longtable head/foot
\IfFileExists{footnotehyper.sty}{\usepackage{footnotehyper}}{\usepackage{footnote}}
\makesavenoteenv{longtable}
\usepackage{graphicx}
\makeatletter
\def\maxwidth{\ifdim\Gin@nat@width>\linewidth\linewidth\else\Gin@nat@width\fi}
\def\maxheight{\ifdim\Gin@nat@height>\textheight\textheight\else\Gin@nat@height\fi}
\makeatother
% Scale images if necessary, so that they will not overflow the page
% margins by default, and it is still possible to overwrite the defaults
% using explicit options in \includegraphics[width, height, ...]{}
\setkeys{Gin}{width=\maxwidth,height=\maxheight,keepaspectratio}
% Set default figure placement to htbp
\makeatletter
\def\fps@figure{htbp}
\makeatother

\usepackage{amsmath, xparse}
\usepackage{fancyvrb, fvextra}
\usepackage{unicode-math}
\usepackage{svg}
\usepackage{multicol}
\usepackage{listings}
\usepackage{systeme}
\usepackage{bm}
\usepackage{xifthen}
\DefineVerbatimEnvironment{Highlighting}{Verbatim}{breaklines,commandchars=\\\{\}}
\lstset{basicstyle=\ttfamily\footnotesize,breaklines=true}
\newcommand\rowop[1]{\scriptstyle\smash{\xrightarrow[\vphantom{#1}]{\mkern-4mu#1\mkern-4mu}}}
\DeclareDocumentCommand\converttorows%
{>{\SplitList{,}}m}%
{\ProcessList{#1}{\converttorow}}
\NewDocumentCommand{\converttorow}{m}
{\ifthenelse{\isempty{#1}}{}{\rowop{#1}}\\}

\DeclareDocumentCommand \rowops{m}
{\;
\begin{matrix}
\converttorows {#1}
\end{matrix}
\; }
\KOMAoption{captions}{tableheading}
\makeatletter
\makeatother
\makeatletter
\makeatother
\makeatletter
\@ifpackageloaded{caption}{}{\usepackage{caption}}
\AtBeginDocument{%
\ifdefined\contentsname
  \renewcommand*\contentsname{Table of contents}
\else
  \newcommand\contentsname{Table of contents}
\fi
\ifdefined\listfigurename
  \renewcommand*\listfigurename{List of Figures}
\else
  \newcommand\listfigurename{List of Figures}
\fi
\ifdefined\listtablename
  \renewcommand*\listtablename{List of Tables}
\else
  \newcommand\listtablename{List of Tables}
\fi
\ifdefined\figurename
  \renewcommand*\figurename{Figure}
\else
  \newcommand\figurename{Figure}
\fi
\ifdefined\tablename
  \renewcommand*\tablename{Table}
\else
  \newcommand\tablename{Table}
\fi
}
\@ifpackageloaded{float}{}{\usepackage{float}}
\floatstyle{ruled}
\@ifundefined{c@chapter}{\newfloat{codelisting}{h}{lop}}{\newfloat{codelisting}{h}{lop}[chapter]}
\floatname{codelisting}{Listing}
\newcommand*\listoflistings{\listof{codelisting}{List of Listings}}
\makeatother
\makeatletter
\@ifpackageloaded{caption}{}{\usepackage{caption}}
\@ifpackageloaded{subcaption}{}{\usepackage{subcaption}}
\makeatother
\makeatletter
\@ifpackageloaded{tcolorbox}{}{\usepackage[skins,breakable]{tcolorbox}}
\makeatother
\makeatletter
\@ifundefined{shadecolor}{\definecolor{shadecolor}{rgb}{.97, .97, .97}}
\makeatother
\makeatletter
\makeatother
\makeatletter
\makeatother
\ifLuaTeX
  \usepackage{selnolig}  % disable illegal ligatures
\fi
\IfFileExists{bookmark.sty}{\usepackage{bookmark}}{\usepackage{hyperref}}
\IfFileExists{xurl.sty}{\usepackage{xurl}}{} % add URL line breaks if available
\urlstyle{same} % disable monospaced font for URLs
\hypersetup{
  colorlinks=true,
  linkcolor={blue},
  filecolor={Maroon},
  citecolor={Blue},
  urlcolor={Blue},
  pdfcreator={LaTeX via pandoc}}

\author{}
\date{}

\begin{document}
\begin{titlepage}

    \newcommand{\HRule}{\rule{\linewidth}{0.5mm}}
    
    \center
    
    \vspace{10cm}

    \textsc{\LARGE Gwinnett School of Math, Science, and Technology }\\[0.3cm]
    
    \vspace{0.5cm}

    \HRule \\[0.4cm]
    { \huge \bfseries Multivariable Calculus Unit 4 Test Project}\\[0.03cm]
    \HRule \\[1.5cm]
    
    \begin{minipage}{0.4\textwidth}
    \begin{flushleft} \Large
    Anish Goyal \\1st Period
    \end{flushleft}
    \end{minipage}
    ~
    \begin{minipage}{0.4\textwidth}
    \begin{flushright} \Large
    Donny Thurston\\Educator
    \end{flushright}
    \end{minipage}\\[1cm]
    
    {\Large 6 November 2023}\\[1cm]
    
    \includegraphics{img/logo.jpeg}\\
    \vfill
\end{titlepage}
\newpage

\ifdefined\Shaded\renewenvironment{Shaded}{\begin{tcolorbox}[enhanced, borderline west={3pt}{0pt}{shadecolor}, interior hidden, breakable, boxrule=0pt, frame hidden, sharp corners]}{\end{tcolorbox}}\fi

\renewcommand*\contentsname{Table of Contents}
{
\hypersetup{linkcolor=}
\setcounter{tocdepth}{4}
\tableofcontents
}
\newpage{}

\hypertarget{does-the-following-set-of-column-vectors-form-a-basis-for-r3}{%
\section{\texorpdfstring{Does the following set of column vectors form a
basis for
\(R^3\)?:}{Does the following set of column vectors form a basis for R\^{}3?:}}\label{does-the-following-set-of-column-vectors-form-a-basis-for-r3}}

\(\symbf{\left\{ \ \ \begin{bmatrix}1 \\ 0 \\ 2\end{bmatrix}, \begin{bmatrix}3 \\ 2 \\ -4\end{bmatrix}, \begin{bmatrix}-3 \\ -5 \\ 1\end{bmatrix} \ \ \right\}}\)

Point out \textbf{\emph{each}} of the requirements for whether a set is
a basis and how these three vectors meet or do not meet each of them.

If it succeeds to form a basis, tell me how it fulfills each
requirement. If it fails, tell me how it fails and succeeds for each
requirement.

\hypertarget{requirement-1-the-set-of-vectors-are-linearly-independent}{%
\subsection{Requirement 1: The set of vectors are linearly
independent}\label{requirement-1-the-set-of-vectors-are-linearly-independent}}

Let's put the vectors into a matrix and check the determinant:
\begin{align*}
&\det\left(\begin{bmatrix}1 & 3 & -3 \\ 0 & 2 & -5 \\ 2 & -4 & 1 \end{bmatrix}\right) \\
&= \det\left(\begin{bmatrix}2 & -5 \\ -4 & 1\end{bmatrix}\right)-3\det\left(\begin{bmatrix}0 & -5 \\ 2 & 1\end{bmatrix}\right) + (-3)\det\left(\begin{bmatrix}0 & 2 \\ 2 & -4\end{bmatrix}\right) \\
&= (2\cdot 20)-3(0\cdot -10)-3(0\cdot -4) \\
&= 40-3(0)-3(0) \\
&= 40
\end{align*} Since \(40 \ne 0\), the vectors are \textbf{linearly
independent}.

\hypertarget{requirement-2-the-set-of-vectors-span-r3}{%
\subsection{\texorpdfstring{Requirement 2: The set of vectors span
\(R^3\)}{Requirement 2: The set of vectors span R\^{}3}}\label{requirement-2-the-set-of-vectors-span-r3}}

Let's check the rank of the previous matrix: \begin{align*}
\mathrm{rref}\left(\begin{bmatrix}1 & 3 & -3 \\ 0 & 2 & -5 \\ 2 & -4 & 1 \end{bmatrix}\right) &= \begin{bmatrix}1 & 0 & 0 \\ 0 & 1 & 0 \\ 0 & 0 & 1 \end{bmatrix}
\end{align*} Since the rank is 3, the vectors \textbf{span
\(\symbf{R^3}\)}.

Therefore, the set of vectors form a \textbf{basis for \(\symbf{R^3}\)}.

\newpage{}

\hypertarget{symbfa-is-a-non-zero-symbf11-times-7-matrix}{%
\section{\texorpdfstring{\(\symbf{A}\) is a non-zero
\(\symbf{11 \times 7}\)
matrix}{\textbackslash symbf\{A\} is a non-zero \textbackslash symbf\{11 \textbackslash times 7\} matrix}}\label{symbfa-is-a-non-zero-symbf11-times-7-matrix}}

What is the largest possible \textbf{\emph{rank}} of \(A\)? What is the
smallest? \textbf{Why?}

The \textbf{largest possible rank} of \(A\) is
\(\mathrm{min}⁡(11,7) = \symbf{7}\) because the minimum of the number of
rows and columns determines the maximum number of linearly independent
rows. The \textbf{smallest possible rank} of \(A\) is \textbf{1}, which
occurs when all the rows/columns of \(A\) are multiples of each other.

What is the largest possible \textbf{\emph{nullity}} of \(A\)?
\textbf{Why?}

The largest possible nullity of \(A\) is the maximum number of linearly
independent columns that can be added to \(A\) before it becomes
singular. This gives us the following relationship:
\(\text{maximum nullity}+\text{minimum rank}=\text{number of rows}\).
Because the minimum rank is 1, \textbf{the maximum nullity is 6} because
\(1+6=7\).

\newpage{}

Make me a particular example of the above matrix, \(\symbf{A}\), that
has a rank of \textbf{4}.

\(\begin{bmatrix} 1 & 0 & 0 & 0 & 0 & 0 & 0\\ 0 & 1 & 0 & 0 & 0 & 0 & 0 \\ 0 & 0 & 1 & 0 & 0 & 0 & 0 \\ 0 & 0 & 0 & 1 & 0 & 0 & 0 \\ 0 & 0 & 0 & 0 & 0 & 0 & 0 \\ 0 & 0 & 0 & 0 & 0 & 0 & 0 \\ 0 & 0 & 0 & 0 & 0 & 0 & 0 \\ 0 & 0 & 0 & 0 & 0 & 0 & 0 \\ 0 & 0 & 0 & 0 & 0 & 0 & 0 \\ 0 & 0 & 0 & 0 & 0 & 0 & 0 \\ 0 & 0 & 0 & 0 & 0 & 0 & 0 \end{bmatrix}\)

What is the nullity of your example?

The \textbf{nullity} is \textbf{3}.

\newpage{}

\hypertarget{symbfa-beginbmatrix2-2-1-4--1-6-4-4-6-0-5--8--4--4--10-8--12-28-0-0--12-24--21-60endbmatrix}{%
\section{\texorpdfstring{\(\symbf{A = \begin{bmatrix}2 & 2 & 1 & 4 & -1 & 6 \\ 4 & 4 & 6 & 0 & 5 & -8 \\ -4 & -4 & -10 & 8 & -12 & 28 \\ 0 & 0 & -12 & 24 & -21 & 60\end{bmatrix}}\)}{\textbackslash symbf\{A = \textbackslash begin\{bmatrix\}2 \& 2 \& 1 \& 4 \& -1 \& 6 \textbackslash\textbackslash{} 4 \& 4 \& 6 \& 0 \& 5 \& -8 \textbackslash\textbackslash{} -4 \& -4 \& -10 \& 8 \& -12 \& 28 \textbackslash\textbackslash{} 0 \& 0 \& -12 \& 24 \& -21 \& 60\textbackslash end\{bmatrix\}\}}}\label{symbfa-beginbmatrix2-2-1-4--1-6-4-4-6-0-5--8--4--4--10-8--12-28-0-0--12-24--21-60endbmatrix}}

\begin{align*}
\mathrm{rref}(A) = \begin{bmatrix}2 & 2 & 1 & 4 & -1 & 6 \\ 4 & 4 & 6 & 0 & 5 & -8 \\ -4 & -4 & -10 & 8 & -12 & 28 \\ 0 & 0 & -12 & 24 & -21 & 60 \end{bmatrix} \rowops{\frac{1}{2}R_1,,,} \begin{bmatrix} 1 & 1 & \frac{1}{2} & 2 & -\frac{1}{2} & 3 \\ 4 & 4 & 6 & 0 & 5 & -8 \\ -4 & -4 & -10 & 8 & -12 & 28 \\ 0 & 0 & -12 & 24 & -21 & 60 \end{bmatrix} \rowops{,R_2-4R_1,R_3+4R_1,} \\
\begin{bmatrix} 1 & 1 & \frac{1}{2} & 2 & -\frac{1}{2} & 3 \\ 0 & 0 & 4 & -8 & 7 & -20 \\ 0 & 0 & -8 & 16 & -14 & 40 \\ 0 & 0 & -12 & 24 & -21 & 60 \end{bmatrix} \rowops{,,R_3+2R_2,} \begin{bmatrix} 1 & 1 & \frac{1}{2} & 2 & -\frac{1}{2} & 3 \\ 0 & 0 & 4 & -8 & 7 & -20 \\ 0 & 0 & 0 & 0 & 0 & 0 \\ 0 & 0 & -12 & 24 & -21 & 60 \end{bmatrix} \rowops{,,,R_4+3R_2} \\
\begin{bmatrix} 1 & 1 & \frac{1}{2} & 2 & -\frac{1}{2} & 3 \\ 0 & 0 & 4 & -8 & 7 & -20 \\ 0 & 0 & 0 & 0 & 0  & 0 \\ 0 & 0 & 0 & 0 & 0 & 0 \end{bmatrix} \rowops{,\frac{1}{4}R_2,,}\begin{bmatrix}1 & 1 & \frac{1}{2} & 2 & -\frac{1}{2} & 3 \\ 0 & 0 & 1 & -2 & \frac{7}{4} & -5 \\ 0 & 0 & 0 & 0 & 0 & 0 \\ 0 & 0 & 0 & 0 & 0 & 0 \end{bmatrix} \rowops{R_1-\frac{1}{2}R_2,,,} \begin{bmatrix} 1 & 1 & 0 & 3 & -\frac{11}{8} & \frac{11}{2} \\ 0 & 0 & 1 & -2 & \frac{7}{4} & -5 \\ 0 & 0 & 0 & 0 & 0 & 0 \\ 0 & 0 & 0 & 0 & 0 & 0 \end{bmatrix}
\end{align*}

Find \ul{\textbf{and justify}} each of the following:

\begin{enumerate}
\def\labelenumi{\alph{enumi}.}
\tightlist
\item
  Set of vectors that form a basis for the column space of \(A\)
\end{enumerate}

We know that the column space of \(A\) is equivalent to:

\(\mathrm{span}\left(\begin{bmatrix}1 \\ 0 \\ 0 \\ 0\end{bmatrix},\begin{bmatrix}1 \\ 0 \\ 0 \\ 0\end{bmatrix},\begin{bmatrix}0\\1\\0\\0\end{bmatrix},\begin{bmatrix}3\\-2\\0\\0\end{bmatrix},\begin{bmatrix}-\frac{11}{8}\\ \frac{7}{4}\\0\\0\end{bmatrix},\begin{bmatrix}\frac{11}{2}\\-5\\0\\0\end{bmatrix}\right)\)

So, in order to get a basis for the column space, we just need to
identify which vectors in the span are linearly independent. We can do
this by looking at the rref of \(A\) and identifying which columns have
pivots. The columns with pivots are the linearly independent columns,
and the columns without pivots are linearly dependent. Therefore, the
basis for the column space of \(A\) is:

\(\left\{\begin{bmatrix}1 \\ 0 \\ 0 \\ 0\end{bmatrix},\begin{bmatrix}0 \\ 1 \\ 0 \\ 0\end{bmatrix}\right\}\)

\newpage{}

\begin{enumerate}
\def\labelenumi{\alph{enumi}.}
\setcounter{enumi}{1}
\tightlist
\item
  Set of vectors that form a basis for the row space of \(A\)
\end{enumerate}

We know that the row space of \(A\) is equivalent to:

\(\mathrm{span}\left(\begin{bmatrix}1 \\ 1 \\ \frac{1}{2} \\ 2 \\ -\frac{1}{2} \\ 3\end{bmatrix},\begin{bmatrix}0 \\ 0 \\ 1 \\ -2 \\ \frac{7}{4} \\ -5\end{bmatrix},\begin{bmatrix}0 \\ 0 \\ 0 \\ 0 \\ 0 \\ 0\end{bmatrix} \right)\)

So, in order to get a basis for the row space, we just need to identify
which vectors in the span are linearly independent. We can do this by
looking at the rref of \(A\) and identifying which rows are non-zero.
The non-zero rows are the linearly independent rows, and the zero rows
are linearly dependent. Therefore, the basis for the row space of \(A\)
is:

\(\left\{\begin{bmatrix}1 \\ 1 \\ \frac{1}{2} \\ 2 \\ -\frac{1}{2} \\ 3\end{bmatrix},\begin{bmatrix}0 \\ 0 \\ 1 \\ -2 \\ \frac{7}{4} \\ -5\end{bmatrix}\right\}\)

\begin{enumerate}
\def\labelenumi{\alph{enumi}.}
\setcounter{enumi}{2}
\tightlist
\item
  Set of vectors that form a basis for the null space of \(A\)
\end{enumerate}

We can represent \(A\) as a system of equations:

\systeme*{x_1 = -x_2-3x_4+\frac{11}{8}x_5-\frac{11}{2}x_6, x_2 = x_2, x_3 = 2x_4-\frac{7}{4}x_5+5x_6, x_4=x_4, x_5=x_5, x_6=x_6}

\(\Rightarrow \vec{x} = \begin{bmatrix}x_1 \\ x_2 \\ x_3 \\ x_4 \\ x_5 \\ x_6\end{bmatrix}=\begin{bmatrix}-x_2-3x_4+\frac{11}{8}x_5-\frac{11}{2}x_6 \\ x_2 \\ 2x_4-\frac{7}{4}x_5+5x_6 \\ x_4 \\ x_5 \\ x_6\end{bmatrix}=x_2\begin{bmatrix}-1 \\ 1 \\ 0 \\ 0 \\ 0 \\ 0\end{bmatrix}+x_4\begin{bmatrix}-3 \\ 0 \\ 2 \\ 1 \\ 0 \\ 0\end{bmatrix}+x_5\begin{bmatrix}\frac{11}{8}\\ 0 \\ -\frac{7}{4} \\ 0 \\ 1 \\ 0\end{bmatrix}+x_6\begin{bmatrix}-\frac{11}{2} \\ 0 \\ 5 \\ 0 \\ 0 \\ 1\end{bmatrix}\)

\newpage{}

Therefore, the basis of the null space can be represented with the
following set:

\(\left\{\begin{bmatrix}-1 \\ 1 \\ 0 \\ 0 \\ 0 \\ 0\end{bmatrix}, \begin{bmatrix}-3 \\ 0 \\ 2 \\ 1 \\ 0 \\ 0\end{bmatrix}, \begin{bmatrix}\frac{11}{8}\\ 0 \\ -\frac{7}{4} \\ 0 \\ 1 \\ 0\end{bmatrix}, \begin{bmatrix}-\frac{11}{2} \\ 0 \\ 5 \\ 0 \\ 0 \\ 1\end{bmatrix}\right\}\)

\begin{enumerate}
\def\labelenumi{\alph{enumi}.}
\setcounter{enumi}{3}
\tightlist
\item
  Rank of \(A\)
\end{enumerate}

The rank of \(A\) is just the \textbf{dimension} of the column space
basis of \(A\). There are two elements in the basis set for the column
space of \(A\), so the rank of \(A\) is \textbf{2}.

\begin{enumerate}
\def\labelenumi{\alph{enumi}.}
\setcounter{enumi}{4}
\tightlist
\item
  Nullity of \(A\)
\end{enumerate}

The nullity of \(A\) is just the \textbf{dimension} of the null space
basis of \(A\). There are four elements in the basis set for the null
space of \(A\), so the nullity of \(A\) is \textbf{4}.

\newpage{}

\hypertarget{prove-the-following}{%
\section{Prove the following:}\label{prove-the-following}}

\textbf{If a matrix \(\symbf{A}\) is not square, then either the row
vectors or the column vectors of \(A\) are linearly dependent.}

I'm going to take an unconventional approach to this proof. Consider a
matrix \(A\) with dimensions \(m \times n\) where \(m > n\). Let
\(\symbf{r_1, r_2, ..., r_m}\) represent the row vectors of \(A\). If
any of the row vectors are the zero vector or equal, they are linearly
dependent by default.

In that case, let's assume none of the row vectors are zero vectors or
equivalent to one another. Since \(m \ne n\), at least two row vectors
must have more components than the possible columns of the matrix by the
Pigeonhole Principle. This implies that there exist constants
\(c_1, c_2, ... c_m\) not all zero, such that:

\(c_1\symbf{r_1}+c_2\symbf{r_2}+ ... + c_m\symbf{r_m} = \symbf{0}\)

This is a non-trivial linear combination of the row vectors that equals
the zero vector, making the row vectors \textbf{linearly dependent.}

But that's only for \(m > n\). What about \(m < n\)?

Well, similar to the approach for the row vectors, if any of the column
vectors is a zero vector, they are linearly dependent. Otherwise, if any
two column vectors are equal, they are linearly dependent because one
can be expressed as a multiple of the other. By the Pigeonhole
Principle, at least two column vectors must have more components than
the number of rows. This implies that there exist constants
\(d_1, d_2, ..., d_n\) not all zero, such that:

\(d_1\symbf{c_1}+d_2\symbf{c_2}+ ... +d_n\symbf{c_n} = \symbf{0}\)

Yet again, this is a non-trivial linear combination of the column
vectors that equals the zero vector, making the column vectors linearly
dependent.

Therefore, if \(A\) is not square, at least one set of vectors
(depending on the values of \(m\) and \(n\)) must be linearly dependent.

\newpage{}

\hypertarget{can-you-go-backwards}{%
\section{Can you go backwards?}\label{can-you-go-backwards}}

I have a matrix that has a \ul{\textbf{row space}} given by all vectors
of the form:

\(\langle -4t + 5w - 3s, 2w + s, 3t + 4w \rangle\)

Provide for me a matrix that has this row space.

We can decomponentize the vectors to get:

\(\begin{bmatrix}-4 \\ 0 \\ 3\end{bmatrix}t + \begin{bmatrix}5 \\ 2 \\ 4\end{bmatrix}w + \begin{bmatrix}-3 \\ 1 \\ 0\end{bmatrix}s\)

Since these vectors form the row space, we know that they have to be
linearly independent. Therefore, we can use them as the rows of a matrix
to get:

\(\begin{bmatrix}-4 & 0 & 3 \\ 5 & 2 & 4 \\ -3 & 1 & 0\end{bmatrix}\)

\emph{Bonus}: Provide for me a matrix that has this \ul{\textbf{null
space}}.

\(\begin{bmatrix}-4 \\ 0 \\ 3\end{bmatrix}t + \begin{bmatrix}5 \\ 2 \\ 4\end{bmatrix}w + \begin{bmatrix}-3 \\ 1 \\ 0\end{bmatrix}s \Rightarrow \begin{bmatrix}-4t+5w-3s \\ 2w + s \\ 3t + 4w\end{bmatrix}=\begin{bmatrix}x_1 \\ x_2 \\ x_3\end{bmatrix} =\)
\systeme*[x_1x_2x_3tws]{x_1 + 4t + 5w -3s = 0, x_2 - 2w - s = 0, x_3 - 3t - 4w = 0}

Which gives us the following matrix:

\(\begin{bmatrix}1 & 4 & 5 & - 3\\ 0 & \end{bmatrix}\)

\newpage{}

\newpage{}

\hypertarget{connect-the-following-equivalent-statements-together-with-an-explanation-of-how-one-statement-implies-the-other.}{%
\section{Connect the following equivalent statements together with an
explanation of how one statement implies the
other.}\label{connect-the-following-equivalent-statements-together-with-an-explanation-of-how-one-statement-implies-the-other.}}

\begin{enumerate}
\def\labelenumi{\arabic{enumi}.}
\tightlist
\item
  \(\symbf{A}x = b\) has exactly one solution for every \(n \times 1\)
  matrix \(b\)
\item
  \(\det{(A)} \neq 0\)
\item
  The columns of \(A\) are linearly independent
\item
  The column vectors of \(A\) span \(R^n\)
\item
  \(A\) has nullity 0
\end{enumerate}



\end{document}
